
%% ------------- English version ---------------
\documentclass[english]{sbrt}
\usepackage[english]{babel}
\usepackage[utf8]{inputenc}
\newtheorem{theorem}{Theorem}
%% ---------------------------------------------

\begin{document}

\title{SBrT 2023 manuscript template}

\author{Name1 Surname1 and Name2 Surname2
\thanks{Name1 Surname1, Department1, University1, City1-State1, e-mail: xxxxx@yyyyy.zzzzz.br; Name2 Surname2, Department2, University2, City2-state2, e-mail: xxxxx@yyyyy.zzzzz.br. This work was partially supported by XXXXXXX (XX/XXXXX-X).}%
}

\maketitle

\markboth{XLI BRAZILIAN SYMPOSIUM ON TELECOMMUNICATIONS AND SIGNAL PROCESSING - SBrT 2023, OCTOBER 08--11, 2023, S\~{A}O JOSÉ DOS CAMPOS, SP}{}


\begin{abstract}
This document is an example of how to use the \LaTeX\ style sbrt.cls to prepare a paper for submission to SBrT~2023. The abstract must have at most 100 words.
\end{abstract}
\begin{keywords}
Paper template, \LaTeX, SBrT~2023.
\end{keywords}


\section{Introduction}

The XLI Brazilian Symposium on Telecommunication and Signal Processing (SBrT~2023) is promoted by the Brazilian Telecommunication Society (SBrT) and, this year, is being organized by the Aeronautics Institute of Tecnology (ITA).

\subsection{About SBrT}

SBrT takes place annually, promoting the most relevant national scientific meeting in the telecommunications area, in which topics of paramount importance to the evolution of research and development of this sector are discussed.

\section{Figures and Tables}
The Table \ref{tab:tabela} is only an example \cite{ref2}.
\begin{table}[htb]
\caption{\label{tab:tabela}The \textit{caption} comes before the table.}
\begin{center}
{\tt
\begin{tabular}{|c||c|c|c|}\hline
&title page&odd page&even page\\\hline\hline
onesided&leftTEXT&leftTEXT&leftTEXT\\\hline
twosided&leftTEXT&rightTEXT&leftTEXT\\\hline
\end{tabular}
}
\end{center}
\end{table}

The Figure \ref{fig:figura} is only an example \cite{ref2}.

\begin{figure}[hbt]
\begin{center}
\setlength{\unitlength}{0.0105in}%
\begin{picture}(242,156)(73,660)
\put( 75,660){\framebox(240,150){}} \put(105,741){\vector( 0, 1){
66}} \put(105,675){\vector( 0, 1){ 57}} \put( 96,759){\vector( 1,
0){204}} \put(105,789){\line( 1, 0){ 90}} \put(195,789){\line(
2,-1){ 90}} \put(105,711){\line( 1, 0){ 60}} \put(165,711){\line(
5,-3){ 60}} \put(225,675){\line( 1, 0){ 72}} \put(
96,714){\vector( 1, 0){204}} \put(
99,720){\makebox(0,0)[rb]{\raisebox{0pt}[0pt][0pt]{a}}}
\put(291,747){\makebox(0,0)[lb]{\raisebox{0pt}[0pt][0pt]{ o}}}
\put(291,702){\makebox(0,0)[lb]{\raisebox{0pt}[0pt][0pt]{ o}}}
\put( 99,795){\makebox(0,0)[rb]{\raisebox{0pt}[0pt][0pt]{ $M$}}}
\end{picture}
\end{center}
\caption{\label{fig:figura}This figure is only an example. The
\textit{caption} must come after the figure.}
\end{figure}

\section{Equation}

This is an example of how to include an equation in the text.

\begin{equation}\label{eq:exemplo}
    h(t)=\sum_{n=0}^{N-1} \alpha_n\delta(t-\tau_n),
\end{equation}
where $\alpha_n$ is the $n$th...

Observe that (\ref{eq:exemplo}) is only an example. There are different ways to include equations with multiple lines, like
\begin{equation} \label{eq:exemplo_multiplo}
    \begin{array}{ccl}
        x^2 & = & bx+c\\
        y^2 & = &\beta y+\gamma=0.
    \end{array}
\end{equation}
Eq. (\ref{eq:exemplo_multiplo}) is only an example.

\section{Conclusions}
The final version of the accepted manuscript for publication in the Proceeding of XLI Brazilian Symposium on Telecommunication and Signal Processing (SBrT~2023) must be sent, in PDF format, at most until the date specified in the call for papers. The manuscript format must be A4, double-column, 10pt, lad-single, and have at most 5 pages. The Abstract must have at most 100 words.

\section*{Aknowledgements}
The Coordination of the SBrT~2023 thanks to the Coordination of the previous symposiums promoted by the Brazilian Telecommunication Society, for making available this example.

\begin{thebibliography}{99}
\bibitem{ref1} L. Lamport, \textit{A Document Preparation System: \LaTeX, User's
Guide and Reference Manual}. Addison Wesley Publishing Company,
1986.
\bibitem{ref2} F. C. Silva e J. J. Sousa, ``This reference is only an example," ~\textit{Journal of examples}, v. 5, pp. 52--55, May
1999.
\end{thebibliography}


\appendix
Insert informations about the apendix here

\end{document}
